%%%%%%%%%%%%%%%%%%%%%%%%%%%%%%%%%%%%%%%%%%%%%%%%%%%%%%%%%%%%%%%%%%%%%%%%%%%%%%%%
% MemCatalog paper commands (ACM SIGMOD / acmart sigconf)
%%%%%%%%%%%%%%%%%%%%%%%%%%%%%%%%%%%%%%%%%%%%%%%%%%%%%%%%%%%%%%%%%%%%%%%%%%%%%%%%

\usepackage{amsthm}
\newtheorem{example}{Example}
\newtheorem{proposition}{Proposition}

\usepackage{xcolor}
\usepackage{multirow}
\usepackage{booktabs}
\usepackage{colortbl}
\usepackage{tikz}
\usepackage{subcaption}
\usepackage{enumitem}
\usepackage{xspace}
\usepackage{array}
\usepackage{makecell}
\usepackage{graphicx}
\usepackage{balance}

\usetikzlibrary{shapes.geometric,shapes.symbols,positioning,arrows.meta,fit,backgrounds,calc}

% ===== Math operators =====
\DeclareMathOperator*{\argmax}{arg\,max}

% ===== System name =====
\newcommand{\sys}{MemCatalog\xspace}
\newcommand{\syst}{MemCatalog\xspace}
\newcommand{\nlsql}{{NL2SQL}\xspace}

% ===== Latin abbreviations =====
\newcommand{\eg}{{\em e.g.,}\xspace}
\newcommand{\ie}{{\em i.e.,}\xspace}
\newcommand{\etc}{{\em etc.}\xspace}
\newcommand{\etal}{{\em et al.}\xspace}

% ===== Paragraph titles (ACM run-in heading convention) =====
\newcommand{\stitle}[1]{\par\smallskip\noindent\textbf{#1}}
\newcommand{\etitle}[1]{\par\smallskip\noindent\textit{#1}}
\newcommand{\runin}[1]{\stitle{#1.}}

% ===== Color palette =====
\definecolor{cadmiumgreen}{rgb}{0.0, 0.42, 0.24}   % positive / gains
\definecolor{dropred}{rgb}{0.75, 0.22, 0.17}       % negative / drops
\definecolor{softblue}{HTML}{2C5F8D}               % primary highlight
\definecolor{softorange}{HTML}{D97F1F}             % secondary highlight
\definecolor{shadecolor}{RGB}{240,240,240}         % subtle background
\definecolor{darkgray_x11}{HTML}{2E2E2E}
\definecolor{lightblue}{HTML}{E8F0F8}
\definecolor{lightgreen}{HTML}{E8F5E9}
\definecolor{lightorange}{HTML}{FFF3E0}

% ===== Result-cell helpers =====
\newcommand{\best}[1]{\textbf{#1}}
\newcommand{\second}[1]{\underline{#1}}
\newcommand{\gain}[1]{{\color{cadmiumgreen}\small(+#1)}}
\newcommand{\loss}[1]{{\color{dropred}\small(-#1)}}

% ===== Query-pipeline labels =====
\newcommand{\opPOINT}{\makecell{\textcolor{softblue}{\textsc{Single-Fact}}\\[-1pt]\textcolor{softblue}{\textsc{Retrieval}}}}
\newcommand{\opAGG}{\makecell{\textcolor{softorange}{\textsc{Independent}}\\[-1pt]\textcolor{softorange}{\textsc{Composition}}}}
\newcommand{\opCMP}{\makecell{\textcolor{cadmiumgreen}{\textsc{Dependent}}\\[-1pt]\textcolor{cadmiumgreen}{\textsc{Composition}}}}

% ===== Benchmark name macros =====
\newcommand{\lmev}{{LME-v2}\xspace}
\newcommand{\lmeo}{{LME-o}\xspace}
\newcommand{\lmes}{{LME-s}\xspace}
\newcommand{\lmem}{{LME-m}\xspace}
\newcommand{\locomo}{{LoCoMo}\xspace}
\newcommand{\evomem}{{EvoMem}\xspace}
\newcommand{\membench}{{MemBench}\xspace}
\newcommand{\dynamicmem}{{DynamicMem}\xspace}

% ===== Query-pipeline shorthands =====
\newcommand{\point}{\textsc{Single-Fact Retrieval}\xspace}
\newcommand{\agg}{\textsc{Independent Composition}\xspace}
\newcommand{\cmp}{\textsc{Dependent Composition}\xspace}

% Core memory notation and the three normalized query-plan classes.
\newcommand{\Hh}{\mathcal{H}}
\newcommand{\Ff}{\mathcal{F}}
\newcommand{\Uu}{\mathcal{U}}
\newcommand{\Mm}{\mathcal{M}}
\newcommand{\ret}{\textsc{Lookup}\xspace}
\newcommand{\red}{\textsc{Independent}\xspace}
\newcommand{\comp}{\textsc{Dependent}\xspace}
\DeclareMathOperator*{\argmin}{arg\,min}
% Examples may continue across columns without repeating their titles.
\newenvironment{memoryexample}[1]{%
  \def\FrameCommand{\fboxrule=0.7pt\fboxsep=5pt\fbox}%
  \setlength{\OuterFrameSep}{6pt}%
  \par\begin{framed}\small\parindent=0pt
  \noindent\textbf{#1}\par\nobreak\smallskip
}{\end{framed}\par}

% ===== Global table typography (ACM booktabs style) =====
\renewcommand{\arraystretch}{1.1}
\setlength{\tabcolsep}{4.5pt}

% Group header rows spanning a booktabs table.
\newcommand{\tblgroup}[2]{\multicolumn{#1}{l}{\cellcolor{gray!12}\textit{#2}}}
\newcommand{\ourgroup}[2]{\multicolumn{#1}{l}{\cellcolor{lightblue}\textit{#2}}}

% Allow LaTeX to stretch line ends slightly to avoid marginal overfull boxes.
\setlength{\emergencystretch}{3em}

% ===== Experiment numbering and compact findings =====
\usepackage{framed}
\newcounter{mycounterexp}
\newcommand{\expnum}{\stepcounter{mycounterexp}\arabic{mycounterexp}}
\newcounter{mycounterfinding}
\newcommand{\findingnum}{\stepcounter{mycounterfinding}\arabic{mycounterfinding}}
\definecolor{shadecolor}{RGB}{242,242,242}
\newcommand{\finding}[1]{%
  \par\smallskip\noindent\textbf{Finding \findingnum.} #1\par\smallskip}
